\section{An arithmetic exclusion of the first mixed-color shells}
\label{sec:mixed-arithmetic-addendum}

This section records a general improvement suggested by the
two-sided Gram viewpoint. Its proof uses only the arithmetic of
\(\mathbb Z[\omega]\) and a rank-four determinant, with no support
graph catalogue or difference-matrix classification.

Throughout let \(3\mid n\), \(H=(\omega^{e_{ij}})\in\muT^{n\times n}\),
and \(Q=\sum_{i<j}|(HH^*)_{ij}|^2\). Assign row \(i\) the
\emph{color} \(\sum_j e_{ij}\pmod3\). For two rows, write
\(a,b,c\) for the multiplicities of \(1,\omega,\omega^2\) in
their entrywise quotient. Then \(a+b+c=n\), and
\[
 \left|a+b\omega+c\omega^2\right|^2
 =a^2+b^2+c^2-ab-bc-ca.
\]
Writing \(x=a-b\) and \(y=b-c\), this norm equals
\(x^2+xy+y^2\). Since \(3\mid n\), we have
\(x\equiv y\pmod3\). Consequently the norm is congruent to
\(0\pmod9\) for two rows of equal color, and to
\(3\pmod9\) for rows of different colors.
In particular, a cross-color pair contributes at least \(3\).

\begin{proposition}[Mixed-color inert-prime gap]
\label{prop:mixed-inert-prime-gap}
Suppose \(3\mid n\) and \(n^n\) is not the norm of an
Eisenstein integer. If the rows of
\(H\in\muT^{n\times n}\) have at least two colors, then
\[
 \boxed{\quad Q(H)\ge 3n+15.\quad}
\]
The same conclusion holds if the columns have at least two colors.
\end{proposition}
\begin{proof}
The norm hypothesis implies that \(n\) is odd and that some
prime \(p\equiv2\pmod3\) occurs to odd order in \(n\).
In particular \(n\ge15\) and \(p\ge5\).
For every nonzero Eisenstein integer \(z\),
the \(p\)-adic valuation of \(\operatorname N(z)\) is even.

Let the three row-color multiplicities be \(r_0,r_1,r_2\).
There are \(r_0r_1+r_0r_2+r_1r_2\) cross-color pairs.
Unless, up to permutation, \((r_0,r_1,r_2)=(n-1,1,0)\),
a mixed coloring has at least \(2(n-2)\) such pairs.
Then
\[
 Q\ge6(n-2)\ge3n+15
\]
for \(n\ge15\), as desired.

It remains to consider the coloring \((n-1,1,0)\).
The Gram defect \(E=HH^*-nI\) is a star with a possible
additional weighted graph among its \(n-1\) majority
vertices. The congruences above imply
\[
 Q=3(n-1)+9k,\qquad k\in\mathbb Z_{\ge0}.
\]
If the conclusion fails, \(k=0\) or \(1\).

When \(k=0\), each of the \(n-1\) star edges has
squared modulus \(3\), and all other off-diagonal
entries vanish. The determinant lemma gives
\[
 \det(HH^*)=n^{n-2}(n^2-3n+3).
\]
The final factor is \(3\pmod p\), so the full determinant
has odd \(p\)-valuation, a contradiction.

When \(k=1\), there are precisely two types.
In the first, a single star edge has squared modulus \(12\)
and the other \(n-2\) edges have squared modulus \(3\).
Then
\[
 \det(HH^*)=n^{n-2}(n^2-3n-6),
\]
whose final factor is \(-6\pmod p\). Again its
\(p\)-valuation is odd.

In the second type all star edges have squared modulus
\(3\), with one additional majority--majority edge
of squared modulus \(9\). If this edge joins vertices
\(1,2\), write its Gram entry as \(a\), and write the
hub entries there as \(g_1,g_2\). Compressing the other
\(n-3\) spokes into one direction of squared norm
\(3(n-3)\) gives the nonzero defect characteristic
polynomial
\[
 x^4-(3n+6)x^2-\tau x+27(n-3),\qquad
 \tau=2\operatorname{Re}(a g_2\overline{g_1})\in\mathbb Z.
\]
Therefore
\[
 \det(HH^*)=n^{n-4}
 \bigl[n^4-(3n+6)n^2+\tau n+27(n-3)\bigr].
\]
The bracket is \(-81\pmod p\), again a \(p\)-adic unit.
Since \(n-4\) and \(v_p(n)\) are odd, we obtain the
same contradiction. (All the displayed determinant
expressions are nonzero modulo the appropriate power
of \(p\), so the valuation argument applies.)

Every remaining mixed-color energy is at least
\(3(n-1)+18=3n+15\).
Finally \(\operatorname{tr}((HH^*)^2)
=\operatorname{tr}((H^*H)^2)\), so row and column
energies agree and the same argument applies to columns.
\end{proof}

\begin{remark}
This proposition concerns mixed colorings. A bound on
\emph{all} third-root matrices at these orders also needs
a separate pure-color argument. The proposition does not
by itself prove the conjectured exact minimum
\(Q=105\) at order \(15\).
\end{remark}
